\mysection{実行可能ファイルの識別}

\subsection{Microsoft Visual C++}
\label{MSVC_versions}

MSVCバージョンとインポート可能なDLL：

\small
\begin{center}
\begin{tabular}{ | l | l | l | l | l | }
\hline
\HeaderColor マーケティング版 & 
\HeaderColor 内部版 & 
\HeaderColor CL.EXE版 &
\HeaderColor インポートされるDLL &
\HeaderColor リリース日 \\
\hline
% 4.0, April 1995
% 97 & 5.0 & February 1997
6		&  6.0	& 12.00	& msvcrt.dll	& June 1998		\\
		&	&	& msvcp60.dll	&			\\
\hline
.NET (2002)	&  7.0	& 13.00	& msvcr70.dll	& February 13, 2002	\\
		&	&	& msvcp70.dll	&			\\
\hline
.NET 2003	&  7.1	& 13.10 & msvcr71.dll	& April 24, 2003	\\
		&	&	& msvcp71.dll	&			\\
\hline
2005		&  8.0	& 14.00 & msvcr80.dll	& November 7, 2005	\\
		&	&	& msvcp80.dll	&			\\
\hline
2008		&  9.0	& 15.00 & msvcr90.dll	& November 19, 2007	\\
		&	&	& msvcp90.dll	&			\\
\hline
2010		& 10.0	& 16.00 & msvcr100.dll	& April 12, 2010 	\\
		&	&	& msvcp100.dll	&			\\
\hline
2012		& 11.0	& 17.00 & msvcr110.dll	& September 12, 2012 	\\
		&	&	& msvcp110.dll	&			\\
\hline
2013		& 12.0	& 18.00 & msvcr120.dll	& October 17, 2013 	\\
		&	&	& msvcp120.dll	&			\\
\hline
\end{tabular}
\end{center}
\normalsize

msvcp*.dllは\Cpp{}関連の関数を持つため、これがインポートされている場合、
おそらく\Cpp{}プログラムです。

\subsubsection{名前修飾}

名前は通常\TT{?}シンボルで始まります。

MSVCの\gls{name mangling}についてはこちらで詳しく読むことができます：\myref{namemangling}。

\subsection{GCC}
\myindex{GCC}

*NIXターゲット以外にも、GCCはCygwinとMinGWの形でwin32環境にも存在します。

\subsubsection{名前修飾}

名前は通常\TT{\_Z}シンボルで始まります。

GCCの\gls{name mangling}についてはこちらで詳しく読むことができます：\myref{namemangling}。

\subsubsection{Cygwin}
\myindex{Cygwin}

cygwin1.dllがしばしばインポートされます。

\subsubsection{MinGW}
\myindex{MinGW}

msvcrt.dllがインポートされる可能性があります。

\subsection{Intel Fortran}
\myindex{Fortran}

libifcoremd.dll、libifportmd.dll、libiomp5md.dll（OpenMPサポート）がインポートされる可能性があります。

libifcoremd.dllには\TT{for\_}で始まる多くの関数があり、これは\IT{Fortran}を意味します。

\subsection{Watcom, OpenWatcom}
\myindex{Watcom}
\myindex{OpenWatcom}

\subsubsection{名前修飾}

名前は通常\TT{W}シンボルで始まります。

例えば、引数を持たず\Tvoidを返すクラス「class」のメソッド「method」は次のようにエンコードされます：

\begin{lstlisting}
W?method$_class$n__v
\end{lstlisting}

\subsection{Borland}
\myindex{Borland Delphi}
\myindex{Borland C++Builder}

これはBorland DelphiとC++Builderの\gls{name mangling}の例です：

\lstinputlisting{digging_into_code/identification/borland_mangling.txt}

名前は常に\TT{@}シンボルで始まり、その後にクラス名、メソッド名、メソッド引数の型がエンコードされます。

これらの名前は.exeインポート、.dllエクスポート、デバッグデータなどにあります。

Borland Visual Component Libraries（VCL）は.dllファイルの代わりに.bplファイルに格納されます。例：vcl50.dll、rtl60.dll。

インポートされる可能性のある別のDLL：BORLNDMM.DLL。

\subsubsection{Delphi}

ほぼすべてのDelphi実行可能ファイルは、コードセグメントの開始時に他の型名と一緒に「Boolean」テキスト文字列を持ちます。

これはDelphiプログラムの\TT{CODE}セグメントの非常に典型的な開始部分で、このブロックはwin32 PEファイルヘッダーの直後に来ます：

\lstinputlisting{digging_into_code/identification/delphi.txt}

データセグメント（\TT{DATA}）の最初の4バイトは\TT{00 00 00 00}、\TT{32 13 8B C0}、または\TT{FF FF FF FF}です。%

この情報は、パック/暗号化されたDelphi実行可能ファイルを扱う際に有用です。

\subsection{その他の既知のDLL}

\myindex{OpenMP}
\begin{itemize}
\item vcomp*.dll---MicrosoftのOpenMP実装。
\end{itemize}